\section{Thermoplaste als Zahnradwerkstoff}
\label{sec:thermoplaste}

Thermoplaste bestehen aus langen, ineinander verknäulten Molekülketten. Innerhalb der Ketten wirken die Hauptvalenzkräfte der chemischen Bindung, zwischen den Ketten nur vergleichsweise schwache Nebenvalenzkräfte \cite[S.~29]{Sto18}. Eine Zugbelastung oder eine Temperaturerhöhung kann diese Verbindung lösen, die Ketten gleiten aneinander ab, und der Werkstoff verformt sich bleibend. Der Übergang zur freien Beweglichkeit der Ketten liegt in einem Temperaturbereich, den die Glasübergangstemperatur kennzeichnet. Unterhalb davon verhält sich ein Thermoplast hart und spröde, oberhalb wird er zunächst formbar und bei weiter steigender Temperatur schmelzeförmig \cite[S.~29]{Sto18}. Beim Abkühlen aus der Schmelze lagern sich die Ketten mancher Polymere teilweise zu kristallinen Überstrukturen zusammen. Diese Thermoplaste heißen teilkristallin, Polymere ohne solche Strukturen amorph \cite[S.~30]{Sto18}. Teilkristalline Thermoplaste bilden ein Zweiphasensystem aus kristallinen Bereichen und amorphen Anteilen an den Lamellenflächen und zwischen den Sphärolithen. Polyamid erreicht bei der Verarbeitung Kristallisationsgrade von $35$ bis $45$\,\%, Polyoxymethylen von $70$ bis $80$\,\% \cite[S.~59 bis 60]{Erh08}.

Amorphe Thermoplaste werden unterhalb ihrer Glasübergangstemperatur eingesetzt und verhalten sich dort hart-spröde. Teilkristalline Thermoplaste arbeiten zwischen der Glasübergangstemperatur ihrer amorphen Bereiche und der Kristallitschmelztemperatur. Die kristallinen Bereiche liefern die Steifigkeit, die erweichten amorphen Bereiche das duktile Verhalten \cite[S.~30 bis 31]{Sto18}. Nach Erhard ist dieser Temperaturbereich der konstruktiv genutzte, weil dort bei ausreichender Steifigkeit die erwartete hohe Zähigkeit vorliegt \cite[S.~70]{Erh08}. Für Zahnräder werden bevorzugt die in der Kurzübersicht der Richtlinie VDI 2736 Blatt 1 aufgeführten Kunststoffe verwendet \cite[S.~5]{VDI2736-1}. Die Kurzübersicht nennt ihre Kennwerte unter Normalklima bei $23$\,\textdegree C und $50$\,\% relativer Feuchte, für die Auslegung verweist die Richtlinie auf herstellerspezifische Werte \cite[S.~5 bis 6]{VDI2736-1}. Amorphe Thermoplaste weisen gegenüber teilkristallinen wie Polyamid, Polyoxymethylen und Polybutylenterephthalat eine geringere Verschleißfestigkeit auf \cite[S.~19]{VDI2736-1}.

Neben der Temperatur beeinflussen Art und Dauer der Belastung, die Verarbeitungsbedingungen und die aufgenommene Feuchte die mechanischen Kennwerte eines Thermoplasts \cite[S.~7]{VDI2736-1}. Mit steigender Temperatur nehmen die molekularen Bindungskräfte ab, und die Steifigkeit sinkt. Bei teilkristallinen Thermoplasten gewährleisten die kristallinen Bereiche auch oberhalb der Glasübergangstemperatur den Zusammenhalt \cite[S.~7]{VDI2736-1}. Ein und derselbe Werkstoff kann je nach Temperatur mehrere charakteristische Formen der Spannungs-Dehnungs-Kurve annehmen, so dass alle Festigkeits- und Dimensionierungskennwerte ausgeprägt temperaturabhängig sind \cite[S.~116 bis 117]{Erh08}. Stommel et al. heben hervor, dass Temperatur- und Zeiteinflüsse bei Kunststoffen schon unter üblichen technischen Einsatzbedingungen zu berücksichtigen sind \cite[S.~31]{Sto18}.

Mit zunehmender Belastungsgeschwindigkeit bleibt den Ketten weniger Zeit zum Entschlaufen. Die viskosen Anteile treten zurück, der Werkstoff reagiert steifer und versagt spröder \cite[S.~30]{Sto18}. Elastizitätsmodul und Streckspannung nehmen bei hohen Belastungsgeschwindigkeiten ebenso zu wie die Zugfestigkeit, die Bruchdehnung nimmt ab \cite[S.~7]{VDI2736-1}. Erhöhte Belastungsgeschwindigkeit und verminderte Temperatur wirken in dieselbe Richtung und lassen sich über das Zeit-Temperatur-Verschiebungsprinzip ineinander überführen \cite[S.~30 und 44]{Sto18}. Der Einfluss der Temperatur auf Festigkeit und Modul ist jedoch deutlich größer als der der Dehngeschwindigkeit \cite[S.~118]{Erh08}. Die Richtlinie stuft ihn ebenso ein, verlangt aber Kennwerte, die zur Art der Belastung sowie zu ihrer Höhe und Dauer passen \cite[S.~7]{VDI2736-1}. Unter konstanter Last kriecht ein Kunststoff, unter konstanter Verformung relaxiert er, und bei dynamischer Beanspruchung sinkt die Festigkeit mit der Lastwechselzahl. Hohe Lastwechselzahlen oder große Dehnungen erwärmen den Werkstoff über sein viskoelastisches Verhalten und mindern die Eigenschaften zusätzlich \cite[S.~7]{VDI2736-1}. Nach ISO 527 wird der Ursprungsmodul bei $1$\,mm/min bestimmt, eine vollständige Kurve häufig bei $50$\,mm/min, so dass der Einpunktwert einer Datenbank vom Ursprungsmodul der Kurve abweicht \cite[S.~31]{Sto18}.

Polyamide nehmen Feuchte aus ihrer Umgebung auf. Die Richtlinie unterscheidet den trockenen, den feuchten und den nassen Zustand, die dem Feuchtegleichgewicht im Vakuum, im Normalklima und im Wasser entsprechen. Feuchte verringert Streckspannung und Elastizitätsmodul, erhöht Zähigkeit und Bruchdehnung und lässt das Volumen wachsen \cite[S.~7 bis 8]{VDI2736-1}. Wassermoleküle lagern sich an die Wasserstoffbrücken der Amidgruppen an und lockern den Zusammenhalt der Ketten. Dadurch verschiebt sich die Glasübergangstemperatur eines Polyamids von etwa $70$\,\textdegree C im trockenen über $10$\,\textdegree C im luftfeuchten auf $-20$\,\textdegree C im nassen Zustand \cite[S.~45]{Erh08}. Bei glasfaserverstärkten Typen hängt die Längenänderung durch Feuchte vom Fasergehalt und von der Faserorientierung ab. In Faserrichtung bleibt etwa ein Viertel der Längenzunahme des unverstärkten Typs \cite[S.~45]{Erh08}. Für Prüfräder aus PA46 beschreibt Hasl eine Nachbehandlung aus Tempern, Konditionieren im Wasserbad und erneutem Tempern, ohne die schon im ungelaufenen Zustand Risse entstehen können \cite[S.~66 bis 67]{Has18}.

Zur Erhöhung von Steifigkeit und Festigkeit werden Thermoplaste mit Kurzglasfasern verstärkt, die Längen von etwa $0{,}2$\,mm und Durchmesser von etwa $10$\,\textmu m besitzen \cite[S.~55]{Erh08}. Eine Verstärkung liegt vor, wenn Festigkeit und Modul des Verbunds über denen der unverstärkten Matrix liegen \cite[S.~55]{Erh08}. Das mechanische Verhalten hängt von der Orientierung der Fasern zur Belastungsrichtung ab. In Faserrichtung werden die höchsten Steifigkeiten und Festigkeiten erreicht, quer dazu deutlich niedrigere \cite[S.~39]{Sto18}. Bei vollständiger Ausrichtung liegt die Festigkeit quer zur Faser sogar unter der der unverstärkten Matrix, weil die Matrix zwischen den Fasern ein Vielfaches der äußeren Dehnung übernehmen muss \cite[S.~66 bis 68]{Erh08}. Auch die Querkontraktionszahl wird richtungsabhängig. Für PA66-GF35 nennt Erhard bei $2$\,mm Wanddicke $0{,}5$ in Spritzrichtung und $0{,}23$ quer dazu \cite[S.~136]{Erh08}. Datenbankwerte faserverstärkter Thermoplaste stammen in der Regel von Probekörpern, in denen verarbeitungsbedingt etwa $80$\,\% der Kurzfasern in Belastungsrichtung liegen \cite[S.~39]{Sto18}. Glasfasern verringern die thermische Längenänderung \cite[S.~9]{VDI2736-1}. Ihre Wirkung auf Schwindung und Verzug behandelt \secref{sec:spritzgiessen}. Illenberger zeigt, dass die höhere Steifigkeit faserverstärkter Werkstoffe die Hertzsche Pressung im Zahnkontakt deutlich erhöht und die Kontaktbreite verringert. Eine Faserverstärkung kann deshalb die Zahnfußtragfähigkeit heben, die Grübchen- und Verschleißtragfähigkeit aber mindern \cite[S.~30]{Ill22}.

Polyamid 46 ist aus Bausteinen mit vier und sechs Kohlenstoffatomen zwischen den Amidgruppen aufgebaut \cite[S.~45]{Erh08}. Es ist teilkristallin, hat im trockenen Zustand eine Glasübergangstemperatur von $75$\,\textdegree C und einen Kristallitschmelzpunkt von $295$\,\textdegree C und nimmt relativ viel Wasser auf \cite[S.~78]{Erh08}. Gegenüber konventionellen Polyamiden wie PA6 oder PA66 zeichnet es sich durch höhere Temperaturbeständigkeit mit Dauereinsatztemperaturen bis etwa $160$\,\textdegree C aus \cite[S.~66]{Has18}. Die Prüfräder dieser Arbeit bestehen aus Stanyl TW200F6 von Envalior, einem mit $30$\,\% Glasfasern verstärkten, wärmestabilisierten PA46 \cite{Env25}. Das CAMPUS-Datenblatt liefert die Kennwerte in \tabref{tab:kennwerte_tw200f6} \cite{Env24}. In \tabref{tab:karten_tw200f6} stehen die Kennwerte der Werkstoffkarten aus den Zugversuchen des Lehrstuhls für Maschinenelemente, Getriebe und Tribologie der Rheinland-Pfälzischen Technischen Universität Kaiserslautern-Landau am konditionierten Werkstoff (MEGT-Zugversuche, Matscape-Karte). Das Datenblatt nennt für den konditionierten Zustand einen Zugmodul von $6000$\,MPa und eine Bruchspannung von $115$\,MPa. Die quasistatische Karte liefert bei $23$\,\textdegree C in Fließrichtung einen Ursprungsmodul von $5980$\,MPa und eine größte Spannung von $115{,}5$\,MPa, quer dazu $3933$\,MPa und $70{,}9$\,MPa. Bei $80$\,\textdegree C fallen die Werte in Fließrichtung auf $3886$\,MPa und $89{,}0$\,MPa. Weil der konditionierte Werkstoff keinen ausgeprägten Streckpunkt zeigt, dient in der Bemessung nach Stommel et al. das Maximum der Kurve in Fließrichtung bei $23$\,\textdegree C als statischer Werkstoffkennwert $K$. Das Maximum bei $80$\,\textdegree C vertritt die Streckspannung $\sigma_{\mathrm{Y}}$ der Betriebsisotherme (vgl. \secref{sec:auswertung_beanspruchung}). Die Karten der beiden höheren Prüfgeschwindigkeiten zeigen den in \secref{sec:materialmodelle} beschriebenen Anstieg von Modul und Festigkeit mit der Dehnrate. Der Abschlussbericht FVA 856 I untersuchte Prüfräder aus PA46-GF30 und entnahm Dichte und Schmelztemperatur sowie die thermischen Kennwerte dem Datenblatt des trockenen Werkstoffs \cite[S.~105 bis 106]{FVA856-I}.

\begin{table}[htbp]
    \caption{Kennwerte des PA46-GF30 Stanyl TW200F6 nach dem CAMPUS-Datenblatt \cite{Env24}. Dabei steht tr.\ für trocken und kond.\ für konditioniert, Striche stehen für nicht angegebene Werte.}
    \label{tab:kennwerte_tw200f6}
    \centering
    \footnotesize
    \setlength{\tabcolsep}{4pt}
    \begin{tabular}{L{5.2cm}L{1.8cm}L{1.8cm}L{1.8cm}L{3.6cm}}
        Kennwert                                        & Einheit                 & tr.                         & kond.          & Prüfnorm       \\
        \hline
        \arrayrulecolor{black!30}
        Dichte                                          & kg/m\textsuperscript{3} & $1410$                      & --             & ISO 1183       \\
        \hline
        Zugmodul                                        & MPa                     & $10\,000$                   & $6000$         & ISO 527-1/-2   \\
        \hline
        Bruchspannung                                   & MPa                     & $210$                       & $115$          & ISO 527-1/-2   \\
        \hline
        Bruchdehnung                                    & \%                      & $3{,}7$                     & $6$            & ISO 527-1/-2   \\
        \hline
        Charpy-Schlagzähigkeit, $23$\,\textdegree C     & kJ/m\textsuperscript{2} & $80$                        & $100$          & ISO 179/1eU    \\
        \hline
        Charpy-Kerbschlagzähigkeit, $23$\,\textdegree C & kJ/m\textsuperscript{2} & $12$                        & $21$           & ISO 179/1eA    \\
        \hline
        Schmelztemperatur, $10$\,K/min                  & \textdegree C           & $295$                       & --             & ISO 11357-1/-3 \\
        \hline
        Glasübergangstemperatur, $10$\,K/min            & \textdegree C           & $75$                        & --             & ISO 11357-1/-2 \\
        \hline
        Formbeständigkeitstemperatur, $1{,}80$\,MPa     & \textdegree C           & $290$                       & --             & ISO 75-1/-2    \\
        \hline
        Feuchtigkeitsaufnahme (Normalklima)               & \%                      & \multicolumn{2}{l}{$2{,}6$} & ähnlich ISO 62                  \\
        \hline
        Wasseraufnahme                                  & \%                      & \multicolumn{2}{l}{$9{,}5$} & ähnlich ISO 62                  \\
        \hline
        \arrayrulecolor{black}
    \end{tabular}
\end{table}

\begin{table}[htbp]
    \caption{Kennwerte der Werkstoffkarten des konditionierten PA46-GF30 je Prüfgeschwindigkeit und Temperatur: Ursprungsmodul und größte Spannung der Zugprobe in Fließrichtung ($0$\textdegree) und quer dazu ($90$\textdegree) aus den MEGT-Zugversuchen (Matscape-Karte).}
    \label{tab:karten_tw200f6}
    \centering
    \footnotesize
    \setlength{\tabcolsep}{4pt}
    \begin{tabular}{L{1.2cm}L{2.6cm}L{1.8cm}L{2.0cm}L{2.0cm}L{2.2cm}L{2.2cm}}
        Karte & Prüfge\-schwindig\-keit in m/s & Temperatur in \textdegree C & Modul $0$\textdegree{} in MPa & Modul $90$\textdegree{} in MPa & größte Spannung $0$\textdegree{} in MPa & größte Spannung $90$\textdegree{} in MPa \\
        \hline
        \arrayrulecolor{black!30}
        QS    & $1{,}67 \cdot 10^{-5}$         & $23$                        & $5980$                        & $3933$                         & $115{,}5$                               & $70{,}9$                                 \\
        \hline
        QS    & $1{,}67 \cdot 10^{-5}$         & $80$                        & $3886$                        & $1787$                         & $89{,}0$                                & $52{,}1$                                 \\
        \hline
        DY1   & $0{,}174$                      & $23$                        & $7590$                        & $5390$                         & $176{,}4$                               & $104{,}9$                                \\
        \hline
        DY1   & $0{,}174$                      & $80$                        & $4626$                        & $2214$                         & $108{,}4$                               & $62{,}5$                                 \\
        \hline
        DY2   & $10$                           & $23$                        & $10\,596$                     & $7923$                         & $245{,}7$                               & $135{,}7$                                \\
        \hline
        DY2   & $10$                           & $80$                        & $6075$                        & $3290$                         & $129{,}8$                               & $75{,}5$                                 \\
        \hline
        \arrayrulecolor{black}
    \end{tabular}
\end{table}

Als Vergleichswerkstoff dient in dieser Arbeit unverstärktes POM, an dem Laufversuche mit derselben Prüfverzahnung aus dem Jahr 2023 vorliegen (vgl. \secref{sec:pruefstand}). Fiederling untersuchte 2023 Räder aus dem Homopolymer Delrin 100NC010 von DuPont und aus PEEK, für das ihre Arbeit sowohl Vestakeep 4000G als auch Vestakeep 4000 GF30 von Evonik nennt \cite[S.~43 bis 44]{Fie23}. Die PEEK-Läufe bleiben außerhalb des Vergleichs. Polyoxymethylen ist ein teilkristalliner Thermoplast mit einer Glasübergangstemperatur um $-60$\,\textdegree C und einem Kristallitschmelzpunkt bei $165$ bis $175$\,\textdegree C \cite[S.~79]{Erh08}. Das homopolymere POM weist gegenüber dem Copolymer tendenziell höhere Festigkeit und Steifigkeit auf und dient wegen seiner Verbreitung und der seit langem verfügbaren Zahnfußfestigkeitswerte als Referenzwerkstoff \cite[S.~66]{Has18}. Polyetheretherketon ist ein teilkristalliner technischer Thermoplast, der zu den Hochleistungskunststoffen zählt und hohe chemische Beständigkeit und Wärmeformbeständigkeit aufweist \cite[S.~43]{Fie23}. Seine Glasübergangstemperatur liegt bei $143$\,\textdegree C, der Kristallitschmelzpunkt bei $344$\,\textdegree C \cite[S.~82]{Erh08}. Gegenüber POM und PA bietet PEEK höhere Festigkeit und Steifigkeit sowie Einsatztemperaturen bis etwa $250$\,\textdegree C bei etwa zwanzigfachem Preis \cite[S.~67]{Has18}. Nach den Herstellerangaben in der Arbeit von Fiederling liegen Elastizitätsmodul und Querkontraktionszahl bei $23$\,\textdegree C für POM bei $2550$\,MPa und $0{,}35$, für PEEK bei $3780$\,MPa und etwa $0{,}4$. PEEK behält seine Steifigkeit bis etwa $120$\,\textdegree C annähernd, während der Modul von POM mit der Temperatur deutlich sinkt \cite[S.~44 bis 45]{Fie23}.

%\Note{PH: Antwort der Betreuerin, welche PEEK-Type 2023 lief (Fie23 nennt 4000G und 4000 GF30); bis dahin kein PEEK-Vergleich.}

% TK erledigt (nur POM 2023 als Vergleich, PEEK-Note bleibt bis zur Antwort): Hier nochmal überprüfen und auch auf die Antwort der Betreuerin abwarten, welche Werkstoffsorten tatsächlich auf dem Prüfstand wann gelaufen sind, da PEEk nur dokumentiert ist, Fiedler jedoch unverstärkt und verstärkt ind er Arbeit ausweist.

% Herkunft:
% Absatz 1 (Aufbau, Glasübergang, amorph/teilkristallin): Sto18 S. 29 (Haupt-/Nebenvalenzkräfte, Abgleiten,
%   Glasübergangstemperatur, hart-spröde, formbar, schmelzeförmig), S. 30 (Kristallisation, teilkristallin,
%   amorph); Erh08 S. 59 (Zweiphasensystem) und S. 60 (Kristallisationsgrade PA 35 bis 45 %, POM 70 bis 80 %).
% Absatz 2 (Einsatzbereich, Zahnradwerkstoffe): Sto18 S. 30 bis 31 (Kurzzeit-Verhalten amorph/teilkristallin,
%   Rolle der Phasen); Erh08 S. 70 (Bereich zwischen Tg und Tm konstruktiv genutzt); Ehr07 S. 220 (7.3.1
%   Werkstoffe, Liste teilkristalliner Thermoplaste, gerenderte Scanseite PDF 98); VDI2736-1 S. 19 (amorph
%   geringere Verschleißfestigkeit), S. 5 bis 6 (Kurzübersicht, Normalklima, herstellerspezifische Kennwerte).
% Absatz 3 (Temperatur): VDI2736-1 S. 7 (5.1 und 5.1.1); Erh08 S. 116 bis 117 (4.1.4 Temperatur); Sto18 S. 31
%   (Bild 2.15, Temperatur- und Zeiteinflüsse unter üblichen Einsatzbedingungen).
% Absatz 4 (Dehnrate, Dauer): Sto18 S. 30 (Entschlaufen, spröder), S. 44 (2.2.5 ZTV-Prinzip), S. 31 (ISO 527,
%   1 mm/min, 50 mm/min); VDI2736-1 S. 7 (5.1.2 Belastungsart und -dauer, Kriechen, Relaxation, Lastwechsel,
%   Erwärmung, Temperatur größer als Belastungsgeschwindigkeit); Erh08 S. 118 (Beanspruchungsgeschwindigkeit,
%   Temperatureinfluss bedeutend größer).
% Absatz 5 (Feuchte): VDI2736-1 S. 7 bis 8 (5.1.3); Erh08 S. 45 (2.2.1 Wasseraufnahme bei Polyamiden,
%   Wasserstoffbrücken, Tg 70/10/-20 °C, Viertel der Längenzunahme in Faserrichtung); Has18 S. 66 bis 67
%   (PA46 Nachbehandlung Tempern, Wasserbad, Tempern; Rissbildung).
% Absatz 6 (Kurzglasfasern): Erh08 S. 55 (2.4.3 Füllung – Verstärkung, Kurzfaser 0,2 mm, 10 µm, Definition
%   Verstärkung), S. 66 bis 68 (2.5.3.2 Füllstofforientierungen, Bild 2.24, Dehnungsvergrößerung, quer unter
%   Matrix), S. 136 (Tabelle 4.5 PA 66-GF 35, 2 mm: 0,5 / 0,23); Sto18 S. 39 (2.2.2, Orientierung zur
%   Belastungsrichtung, 80 % der Fasern in Probekörpern); VDI2736-1 S. 9 (5.2 Glasfasern verringern thermische
%   Längenänderung), S. 25 (7.1 Schwindung, Verzug durch Faserorientierung); Ehr07 S. 220 (Faserzugabe,
%   Verschleiß); Ill22 S. 30 (Hertzsche Pressung, Kontaktbreite, Zahnfuß gegen Flanke).
% Absatz 7 (PA46, TW200F6, Tabellen): Erh08 S. 45 (Fußnote PA 46: 4 und 6 C-Atome), S. 78 (Kurzcharakteristik
%   PA 46 trocken: teilkristallin, Tg 75 °C, Tm 295 °C, relativ hohe Wasseraufnahme); Has18 S. 66 (PA46
%   Dauereinsatztemperatur bis ca. 160 °C gegenüber PA6, PA66, PA12); Env25 (StanylTW200F6_IMR.txt Zeile 5 bis 7:
%   "PA46-GF30, 30% Glass Reinforced, Heat Stabilized"); Env24 (StanylTW200F6.txt Zeilen 14 bis 35: Dichte,
%   Zug-Modul 10000/6000, Bruchspannung 210/115, Bruchdehnung 3.7/6, Charpy 80/100 und 12/21, Schmelz- und
%   Glasübergangstemperatur, Formbeständigkeit, Feuchtigkeits- und Wasseraufnahme); Matscape-Karte:
%   stommel_materials.yaml model_properties der Karten pa46_gf30_qs/dy1/dy2 (E_0, E_90, sigma_b_0, sigma_b_90
%   bei 23 und 80 °C) und Prüfgeschwindigkeiten aus den designation-Zeilen; K als Maximum und sigma_Y der
%   Isotherme: stommel_materials.yaml (static_characteristic_value maximum, validation_matscape sigma_B_S) und
%   kap3_7_auswertung.tex Absatz 4; FVA856-I S. 105 bis 106 (9.4 Prüfzahnräder PA46-GF30, Tabelle 16 aus dem
%   Datenblatt [CAM19] für trockenen Werkstoff).
% Absatz 8 (POM, PEEK): Fie23 S. 43 bis 44 (4.2.1 Werkstoffe, Delrin 100NC010 DuPont, Vestakeep 4000G Evonik,
%   PEEK Hochleistungskunststoff, Tabelle 4: POM 2.550 MPa, 0,35; PEEK 3.780 MPa, ca. 0,4), S. 45 (Abbildung 29,
%   E-Modul über Temperatur); Erh08 S. 79 (POM Tg -60 °C, Tm 165 bis 175 °C), S. 82 (PEEK Tg 143 °C, Tm 344 °C);
%   Has18 S. 66 (Delrin 100NC010 Homopolymer, Referenzwerkstoff), S. 67 (PEEK Festigkeit, Steifigkeit,
%   Einsatztemperatur ca. 250 °C, 20-facher Preis).
% Seitenzuordnung: Sto18 Druckseite = PDF-Seite - 18 (Kopfzeilen der Textdatei); Erh08 Kopfzeilen der
%   Textdatei (Seitenzahl am Seitenanfang); VDI2736-1 Kopfzeilen "– N –"; Has18 Kopfzeilen; Fie23 und Ill22 mit
%   pdftotext an der PDF geprüft (MA-0296 PDF 56 bis 58 = S. 43 bis 45; Illenberger PDF 44 = S. 30);
%   FVA856-I Marker PAGE 119/120 = S. 105/106.
% Bildwunsch: keiner (Tabellen statt Bild; ein Bild aus einer Quelle ist nicht nötig).
% Bibwunsch: Erh08, Env25 übernommen (litbib.bib); Ehr07 gestrichen, weil der Scan kein Impressum trägt und die
%   Auflage nicht prüfbar ist (Hauptagent 2026-10-18): die beiden Ehr07-Aussagen auf VDI 2736 Blatt 1 S. 5 und
%   S. 19 umgestellt bzw. gestrichen (Faserzugabe/Verschleiß steht mit Illenberger S. 30). Env24, Sto18,
%   VDI2736-1, Has18, Ill22, Fie23, FVA856-I liegen in litbib.bib.
% Nomenklatur: K, MPa, statischer Werkstoffkennwert (steht in Nomenklatur.tex); sigma_Y, MPa, Streckspannung
%   (statisch) (steht in Nomenklatur.tex); keine neuen Zeichen.
% Offen: BA-0486 Abschnitt 2.1 (Spritzgießen, IMS-Zahnrad, Bon14, İmr09, Sta20) nicht übernommen, gehört zu 2.2
%   oder ohne Quelle; BA 2.2.1 Einordnung nach Bau22/Dom12 (Werkstoffklassen, Polymerbegriff, Thermoplast/
%   Elastomer/Duroplast) gestrichen, Quellen fehlen; Bon14 (teilkristallin zäher) auf Sto18 S. 30 bis 31 und
%   Erh08 S. 61 umgestellt; Bon14 (Standard-, technische, Hochleistungsthermoplaste) und Man13 (Pyramide)
%   gestrichen; Tabelle 2.1 der BA (16MnCr5 gegen PEEK 5000G, CWF24c) und Den12 (Wärmeausdehnung) gestrichen,
%   ersetzt durch VDI2736-1 S. 9 (Glasfasern, thermische Längenänderung); Rös05 (unvernetzte Thermoplaste
%   etabliert) auf Ehr07 S. 220 und VDI2736-1 S. 5 umgestellt; Bon14 (Polarität, Schmierstoff), Dom12 und
%   Boh10 (Polykondensation, PA66-Namensgebung, PEEK-Monomere) gestrichen, PA46-Namensgebung nach Erh08
%   S. 45; Ill22 (Faserverstärkung Fuß gegen Flanke) an S. 30 geprüft und übernommen.
%   Ehr07: der Scan unter 04_buecher hat kein Impressum; Auflage und Jahr sind nicht prüfbar (Bibliotheksstempel
%   95 A 2861 und Einleitungsbilder bis 1995/1997 deuten auf eine ältere Auflage als die 3. von 2007); die
%   Seiten 220 (7.3.1 Werkstoffe) und 10 bis 11 (1.1.3) wurden am Scan gelesen und müssen gegen die
%   zitierte Auflage geprüft werden.
%   ErhStr: der Scan "Erhard, Strickle - Maschinenelemente aus thermoplastischen Kunststoffen" (Band 2, Lager und
%   Antriebselemente, VDI-Verlag 1978, ISBN 3-18-400369-8) enthält nur S. 1 bis 109 (Gleitlager, Wälzlager,
%   Laufrollen, Seilrollen); das Kapitel 6 Zahnräder ab S. 127 fehlt, deshalb nicht zitiert.
%   Die Feuchtigkeits- und Wasseraufnahme des Datenblatts stehen dort in der Spalte "tr. / kond." als "2.6 / *"
%   und "9.5 / *"; welcher Zustand gemeint ist, sagt das Blatt nicht (hier als Gleichgewichtswerte gelesen).
%   Die Prüfgeschwindigkeiten der Karten (1,67e-5, 0,174, 10 m/s) stammen aus den Kartennamen in
%   stommel_materials.yaml; die Dehnraten dazu stehen in tab:raten (3.6).
